\emph{Tools.} We use the two elementary tools of the proof of \claimref{synthesis.wall-class-one-rung-obstruction}: (a) a nondegenerate symmetric form on a bundle over a compact base has a maximal negative subbundle $N(h)$, unique up to isomorphism; (b) a continuous family of subbundles along a collar is constant up to isomorphism, so bundles over $C$ are identified with bundles over $\Sig$. From classical bundle theory (Fact \claimref{sheaf.imported-real-line-bundles}, and \cite[\S 4]{MilnorStasheff1974} for the two further classical inputs, cited from memory): the total Stiefel-Whitney class is multiplicative, $w(E \oplus E') = w(E)\,w(E')$, with $w_j(E) = 0$ for $j > \operatorname{rank} E$; $H^*(\R P^n; \mathbb{Z}_2) = \mathbb{Z}_2[x]/(x^{n+1})$ with $x = w_1$ of the tautological line $\gamma$; the Euler sequence $T\R P^n \oplus \underline{\R} \cong (n{+}1)\gamma$; and cellular approximation, by which a map from a $d$-dimensional complex to $\R P^\infty$ deforms into $\R P^d$ (\cite[\S\S 4--5]{MilnorStasheff1974}, from memory). Throughout, $(M,G)$ satisfies (T1) to (T3) (\claimref{sheaf.tameness}), and on the physical half of the collar $N(G) \cong N(\overline G)$, as shown in the proof of \claimref{synthesis.wall-class-one-rung-obstruction}.

\emph{(1).} Let $h$ be a $k$-rung completion and $V := N(h)$ over the physical half $C_+ = \Sig \times [0,1]$ of the collar. At $\Sig \times \{1\}$, $h = G \perp H_k$, so $V \cong N(\overline G) \oplus \underline{\R}^{k}$ there, hence along $\Sig$ by tool (b). At $\Sig$, let $V_1$ be a complement of $\shL$ in $T^*M|_\Sig$; since $h|_{T^*M} = G$ along $\Sig$, $h|_{V_1} = \overline G$ is nondegenerate of signature $(p{-}1, q)$ (\claimref{sheaf.degeneration-sheaf}), so $T_k|_\Sig = V_1 \oplus W$ with $W := V_1^{\perp_h}$, of rank $(n + 2k) - (n - 1) = 2k + 1$ and signature $(p{+}k, q{+}k) - (p{-}1, q) = (k{+}1, k)$. The line $\shL$ lies in $W$ ($h(\shL, T^*M) = G(\shL, \cdot) = 0$ on $\Sig$) and is isotropic. Choose a line subbundle $\shL''$ of $W$ complementary to $\shL^{\perp} \cap W$; $h$ pairs $\shL''$ nondegenerately with $\shL$, so $\shL'' \cong \shL^{*} \cong \shL$, and $P := \shL \oplus \shL''$ is a nondegenerate plane bundle of signature $(1,1)$ (Gram $\begin{psmallmatrix} 0 & a \\ a & b \end{psmallmatrix}$, $a \neq 0$). Then $W = P \perp R$ with $R := P^{\perp} \cap W \cong \shL^\perp/\shL$, nondegenerate of signature $(k, k{-}1)$. The negative line of $P$ meets the isotropic line $\shL$ trivially, so it maps isomorphically onto $P/\shL \cong \shL'' \cong \shL$; hence $N(W) \cong \shL \oplus N(R)$, and $V|_\Sig \cong N(\overline G) \oplus \shL \oplus N$ with $N := N(R)$ of rank $k - 1$. Comparing the two descriptions of $V|_\Sig$ gives the isomorphism in (1). Taking total Stiefel-Whitney classes, $w(N(\overline G))\,(1 + x)\,w(N) = w(N(\overline G))$, and $w(N(\overline G))$ is invertible, so $w(N) = (1+x)^{-1} = \sum_j x^{j}$; since $\operatorname{rank} N = k - 1$, $x^{j} = w_j(N) = 0$ for $j \geq k$. For $k = 1$ this is clause (4) of \claimref{synthesis.wall-class-one-rung-obstruction}. The check verifies the inertia identity $\operatorname{In}(W) = (1,1) + \operatorname{In}(\shL^\perp/\shL)$ for $k = 1, 2, 3$ on rational blocks presented in a mixed basis, and the truncated inverse $(1+x)^{-1}$ in $\mathbb{Z}_2[x]/(x^{n+1})$ for $n \leq 7$.

\emph{(2).} Let $\iota : \shL \hookrightarrow \underline{\R}^{k}$ and $N := \iota(\shL)^{\perp}$ for the standard metric, so $\underline{\R}^{k} = \iota(\shL) \oplus N$ over $\Sig$, hence over $C$ by pullback. Present $H_k$ as $\mathcal V \oplus \mathcal V^{*}$ with $\mathcal V = \underline{\R}^{k}$ null and the duality pairing. Over $C$, with $\Lambda := \iota(\tilde\shL)$ and the dual splitting $\mathcal V^* = \Lambda^{*} \oplus N^{*}$ (annihilators), $H_k = (\Lambda \oplus \Lambda^{*}) \perp (N \oplus N^{*})$, each summand hyperbolic. Take twisted collar data $(u, \beta, \mu, \nu)$ for $\Lambda$ as in \claimref{synthesis.twisted-wall-induction-conserves-dying-line-class} (they need $\Lambda$ only on the support of $\beta$, which lies in $C$, and $\Lambda|_\Sig \cong \shL$ via $\iota$), and set $h := (G \perp H_\Lambda + \kappa) \perp H_N$ on $C$, $h := G \perp H_k$ off the support of $\beta$. The two definitions agree where $\beta = 0$, because the splitting of $H_k$ is an orthogonal decomposition of the same form; so $h$ is a smooth form on the trivial bundle $T_k$, equal to $G \perp H_k$ outside $C$, with $h|_{T^*M} = G$ everywhere since $\kappa$ pairs $T^*M$ only with the rung. In a local unit frame of $\Lambda$, the first summand is the ambient form of \claimref{obj.sheaf.collar-data-coupling} (clause (1) of the twisted result) and the second is a fixed hyperbolic block, so $\det h = \pm \det G \cdot F$ with $F = 1 - 2\mu\nu\beta^{2}s$ and $s = G^{-1}(u,u)$. Take $\mu\nu < 0$ and good couplings in the sense of \claimref{sheaf.tame-wall-induction}, and in addition $2|\mu\nu| S_+ < 1$ with $S_+$ the supremum of $|s|$ over the physical half of the collar outside a neighbourhood of $\Sig$ on which $s > 0$ (finite by compactness, since $s = 1/\lambda + O(1)$ with $\lambda$ the dying eigenvalue). Then $F > 0$ on the physical side, and on $\Sig$ the wall Gram computation of \claimref{sheaf.wall-gram} applies to the first summand, as in \claimref{sheaf.one-sided-completion}; so $h$ is nondegenerate on $\Mp$ including $\Sig$, of signature $(p{+}k, q{+}k)$. On the far side, clauses (1) and (2) of the twisted result apply verbatim inside $C$: the shadow wall $\Sig' = \{F = 0\}$ is a tame corank-one wall, a graph over $\Sig$ in the collar, and its radical projects injectively to $T^*M$ along $\Lambda$, so $\shL' \cong \Lambda|_{\Sig'} \cong \shL$ under the collar identification. No extension of $\shL$ over $\Mp$ is used. The check builds this form on the stored torus wall with $k = 2$ (below) and verifies it.

\emph{(3).} The lower bound is (1); the upper bound is (2); and $e(\shL) \leq d + 1$ for $d = \dim \Sig$, because the classifying map $\Sig \to \R P^{\infty}$ of $\shL$ deforms into $\R P^{d}$ (cellular approximation), so $\shL \cong f^{*}\gamma \subset f^{*}\underline{\R}^{d+1}$. For $k = 1$, $\operatorname{ht}(x) = 1$ iff $x = 0$, and $e = 1$ iff $\shL$ is trivial, iff $x = 0$.

\emph{The torus wall.} On $M = T^{2} \times \R$ with $G = -\partial_\varphi^{2} + t K^{2} + M^{2}$ (the example of \claimref{synthesis.twisted-wall-induction-conserves-dying-line-class}(4), $(c, s) = (\cos\frac\varphi2, \sin\frac\varphi2)$), $\shL = \operatorname{span}(\theta)$, $\theta = c\,d\psi + s\,dt$ (written $\theta$ here to keep $k$ for the number of rungs), and $x \neq 0$; $x^{2} = 0$ since $H^*(T^2; \mathbb{Z}_2)$ is an exterior algebra and $x$ is pulled back from the $\varphi$ circle. So $k(\Sig) \geq 2$ by \claimref{synthesis.wall-class-one-rung-obstruction}. The map $\theta \mapsto l := c\,a_{1} + s\,a_{2}$ embeds $\shL$ in the trivial null plane $\mathcal V = \operatorname{span}(a_1, a_2)$ of $H \perp H$ (basis $a_1, a_2, b_1, b_2$, $H(a_i, b_j) = \delta_{ij}$), with complement $N = \operatorname{span}(n)$, $n = -s\,a_1 + c\,a_2$. With $u = c\,\partial_\psi + s\,\partial_t$ the form of (2) has entries $h(\alpha, a_i) = \beta\mu\, l_i\, \alpha(u)$ and $h(\alpha, b_i) = \beta\nu\, l_i\, \alpha(u)$, quadratic in $(c,s)$, hence single valued on the torus. The check verifies, at five rational points of the half-angle circle, $\beta \in \{1, \frac12, 0\}$, $(\mu,\nu) = (-\frac13, 1)$ and five values of $t$: invariance under $(c,s) \mapsto (-c,-s)$; $h|_{T^*M} = G$; $h = G \perp H \perp H$ where $\beta = 0$; $\det h = 2\mu\nu\beta^{2} - t$ (nonzero for $t \geq 0$, $\beta \neq 0$); signature $(4,3)$ for $t \geq 0$, the wall included; and at $t = 2\mu\nu\beta^{2}$ corank one with radical having cotangent part along $\theta$ and rung part along $(c,s)$ in $\mathcal V$ and in $\mathcal V^*$, so $\shL' \cong \shL$. In the Clifford algebra of $\mathcal V$ (Euclidean, GA) $l \wedge n = a_1 \wedge a_2$, so $\shL \oplus N = \mathcal V$ with a global oriented frame. Hence $k(\Sig) = 2$.

\emph{Projective walls.} On $\Sig = \R P^{n}$, by the Euler sequence $T^*\Sig \oplus \underline\R \cong (n{+}1)\gamma$ contains a copy $\shL$ of $\gamma$. On $M = \Sig \times \R$ with a bundle metric $A$ on $T^*M$ and $\pi$ the $A$-orthogonal projection onto the pullback of $\shL$, the cometric $G := A - (1 - t)\,A(\pi\,\cdot, \pi\,\cdot)$ has eigenvalue $t$ on $\shL$ and $1$ on its complement: corank one exactly at $t = 0$, radical $\shL$, $\det G \propto t$ with a simple zero (so (T3) holds by Jacobi's formula, as in \claimref{sheaf.tame-wall-induction}(3)), signature $(n{+}1, 0)$ for $t > 0$ and $(n, 1)$ for $t < 0$. Then $x^{n} \neq 0$ gives $k(\Sig) \geq n + 1$, and $\gamma \subset \underline{\R}^{n+1}$ tautologically gives $k(\Sig) \leq n+1$. The check verifies the Whitney step: $(1+x)^{-1}$ in $\mathbb{Z}_2[x]/(x^{n+1})$ has $w_n \neq 0$, so the least admissible rank of $N$ is $n$.

\emph{(4).} For quadratic bundles $\Cl(A \perp B) \cong \Cl(A) \hat\otimes \Cl(B)$ (graded tensor product). On the collar $T^*M = Q \oplus \tilde\shL$ with $G = G|_Q \perp \hat\lambda$, $\hat\lambda$ of the side's sign (proof of \claimref{synthesis.wall-class-one-rung-obstruction}). Write $\shL_\pm$ for $\shL$ with a positive, respectively negative, fibre metric. Then, over $\Sig$ via tool (b),
\[
\Cl(G_{\mathrm{far}}) = \Cl(G|_Q) \hat\otimes \Cl(\shL_-) \sim \Cl(G|_Q) \hat\otimes \Cl(\shL_+) \hat\otimes \Cl(\shL_-) \hat\otimes \Cl(\shL_-) = \Cl(G_{\mathrm{phys}}) \hat\otimes \Cl(\shL \otimes \R^{0,2}),
\]
where $\sim$ inserts $\Cl(\shL_+ \perp \shL_-) = \Cl(\shL \otimes \R^{1,1})$, which is $\operatorname{End}$ of a graded bundle: with local generators $l_0 = \hat l \otimes e_0$, $l_1 = \hat l \otimes e_1$ ($l_0^2 = 1$, $l_1^2 = -1$), $f = \frac12(1 + l_0 l_1)$ is a primitive idempotent unchanged by the sign change $\hat l \mapsto -\hat l$, so $\Cl \cdot f$ is a global spinor bundle. On the torus wall, $\hat l \otimes e_1 \mapsto l$, $\hat l \otimes e_2 \mapsto n$ is an isometry $\shL \otimes \R^{0,2} \cong (\mathcal V, -\text{metric})$ (both images flip together under the monodromy), so $\Cl(\shL \otimes \R^{0,2})$ is the constant bundle $\Cl(0,2)$ with global pseudoscalar $a_1 a_2$, $(a_1 a_2)^2 = -1$, exactly as for a wall with trivial dying line. So the jump of the Clifford sheaf across the torus wall equals that across an untwisted wall, while $x \neq 0$ obstructs one rung there; since the jump is a function of $\shL$ alone, $W(\Sig)$ determines it, and the torus wall shows the converse fails. The check verifies $f^2 = f$, $f \neq 0, 1$, its invariance under the joint sign flip, and $(a_1a_2)^2 = -1$, in GA.

\emph{Remark (literature, from memory; not used above, and not proved here beyond the torus wall).} In the graded Brauer group of Donovan and Karoubi, $GBO(\Sig) \cong H^0(\Sig;\mathbb{Z}_8) \times H^1(\Sig;\mathbb{Z}_2) \times H^2(\Sig;\mathbb{Z}_2)$ as a set, the class of $\Cl(\shL \otimes \R^{0,2})$ relative to the constant $\Cl(0,2)$ is $(0, w_1(2\shL), w_2(2\shL)) = (0, 0, x^{2})$, so the jump of the sheaf-level mod-8 class across the wall is the clock step $-2$ together with $x^2$. Reading (argued, resting on that remark): the torsion that the sheaf-level mod-8 class can see would be $x^2$, which is the primary part of the two-rung condition in (1); but $x^{2} = 0$ does not imply that two rungs suffice (next remark), so this is a statement about what the class sees, not about completions. The founding pattern ``moved, never destroyed'' holds for the class itself (carried to every shadow wall by (2)) and for the clock step, but not for the obstructive force of the torsion, which $\operatorname{ht}(x)$ rungs absorb.


\emph{Remark (the gap, argued; classical inputs from memory).} The bounds in (3) can differ. On $\Sig = L(4,1)$, $H^1(\Sig;\mathbb{Z}) = 0$, so the Bockstein $\beta : H^1(\Sig;\mathbb{Z}_2) \to H^2(\Sig;\mathbb{Z}) = \mathbb{Z}_4$ is injective and $\beta x = 2$; $x$ lifts to $\mathbb{Z}_4$, so $x^{2} = Sq^{1}x = 0$ and $\operatorname{ht}(x) = 2$. The Euler class of the oriented plane $\shL \oplus \shL$ is $\beta x \neq 0$, so $\shL \oplus \shL$ is not trivial and $e(\shL) \geq 3$. If $N(\overline G) = 0$ (the physical side positive definite, $q = 0$), the isomorphism in (1) with $k = 2$ reads $\shL \oplus N \cong \underline{\R}^{2}$ with $N \cong \shL$ (determinant lines), which is impossible; so $k(\Sig) \geq 3 > \operatorname{ht}(x)$. Such walls exist: on $L(4,1) \times \R$ the line $\shL$ embeds in $T^*\Sig \oplus \underline{\R}$ (rank $4 > 3$, general position), and the cometric of the projective-wall construction is positive definite for $t > 0$. If $N(\overline G) \neq 0$ (for instance a Lorentzian physical side), the necessary condition of (1) holds at $k = 2$ (over a $3$-complex an oriented bundle is stably trivial iff $w_2 = 0$, and $w_2(2\shL) = x^{2} = 0$), but $\shL \oplus \shL \neq \underline{\R}^{2}$, so the collar-confined construction of (2) fails; whether $k(\Sig) = 2$ there is open.
